% ============================================================
% WS3 — architecture-claims rewrite (drop-in for sn-article.tex)
%
% Basis: completed 2x2 controlled comparison (2D TFI, cross-scale).
%   masked diffusion (lr 1e-5 / 1e-4):  5.6e-4, 1.3e-3, 3.0e-3  /  8.7e-4, 2.6e-3, 3.9e-3
%   AR twin        (lr 1e-4 / 1e-5):    3.8e-4, 4.3e-4, 8.1e-4  /  4.1e-4, 4.6e-4, 9.2e-4
%   (mean MSE over three ZZ observables at 4x4 / 4x5 / 4x6)
% The AR twin is identical in tokens, embeddings, size, RoPE, and data;
% only the attention mask and factorization differ.
%
% WHAT TO REMOVE from the current draft:
%  * Introduction: the claim that causal attention is "less compatible
%    with the nature of quantum many-body systems" as an ACCURACY
%    argument, and "autoregressive inductive biases may introduce an
%    unnecessary sequential bias" — soften both as described below.
%  * Results ("Diffusion generative modeling"): "The choice of full
%    attention over causal attention is not merely an architectural
%    preference but an alignment with the physical principles..." —
%    replace with the paragraph below.
% ============================================================

% ---------------------------------------------------------------
% (1) Replacement paragraph for the Results architecture discussion
% ---------------------------------------------------------------

To isolate the role of the generative factorization, we trained a
controlled autoregressive counterpart of Diffushadow: identical token
layout, embeddings, network size, positional encoding, and training data,
differing only in a causal attention mask and a fixed left-to-right
factorization of the measurement distribution. On the two-dimensional
cross-scale benchmark the two models achieve comparable accuracy, with the
autoregressive variant performing at least as well on the site-averaged
correlations studied here (Extended Data Table~X). This is an honest and,
we believe, informative outcome: for ground states of stoquastic
Hamiltonians measured in product bases, a fixed-order factorization of the
measurement distribution is not an obstacle to learning size-transferable
structure. The advantages of the masked-diffusion formulation lie
elsewhere, in capabilities rather than raw accuracy on fixed observables.
First, a single trained model conditions on \emph{any} subset of observed
outcomes: it can complete partially observed snapshots (measurement
inpainting) for arbitrary patterns of known and unknown sites, which a
fixed-order autoregressive model supports only when the conditioning set
is a prefix of its generation order. Second, the denoising order is chosen
adaptively at generation time. Our decoding analysis shows that this
scheduling freedom carries physical weight: unmasking many sites per step
samples them conditionally independently and destroys long-range
correlations in ordered phases, while entropy-guided sequential decoding
--- in which the model itself selects which site to commit next ---
restores them (Extended Data Fig.~Y). The masked-diffusion framework
therefore provides a strictly larger family of conditioning and decoding
strategies within one model, and the attention analysis of the following
section shows that its bidirectional attention organizes information in
patterns that mirror physical correlation structure.

% ---------------------------------------------------------------
% (2) Introduction: replace the last two sentences of the
% autoregressive-vs-diffusion paragraph with:
% ---------------------------------------------------------------

While autoregressive models impose a fixed generation order, masked
diffusion models the joint measurement distribution under arbitrary
conditioning patterns and chooses its generation order adaptively. We
compare the two factorizations in a controlled setting below; we find
that the practical strengths of the diffusion formulation lie in this
conditioning and scheduling flexibility, together with bidirectional
attention structures that align with physical correlations, rather than
in a generic accuracy advantage.

% ---------------------------------------------------------------
% (3) Extended Data Table X (the 2x2 grid) — caption
% ---------------------------------------------------------------

% Extended Data Table X: Controlled factorization comparison on the
% two-dimensional transverse-field Ising benchmark. Mean squared error of
% three site-averaged ZZ observables estimated from 20,000 generated
% snapshots per coupling, at the trained size (4x4) and two unseen sizes.
% The autoregressive (AR) twin shares tokens, embeddings, network size,
% positional encoding and training data with Diffushadow; only the causal
% attention mask and the fixed generation order differ.  Both models are
% shown at two learning rates; results are robust to this choice.
%
%                       4x4 (trained)   4x5 (unseen)   4x6 (unseen)
%  Diffushadow (lr 1e-5)   5.6e-4          1.3e-3         3.0e-3
%  Diffushadow (lr 1e-4)   8.7e-4          2.6e-3         3.9e-3
%  AR twin     (lr 1e-4)   3.8e-4          4.3e-4         8.1e-4
%  AR twin     (lr 1e-5)   4.1e-4          4.6e-4         9.2e-4

% ---------------------------------------------------------------
% (3b) Measurement-inpainting demonstration (Extended Data figure/table)
% Data: results2d/inpaint_demo.npz (4x4 TFI torus, 2000 exact ED test
% snapshots per coupling, hidden-half completions).
% ---------------------------------------------------------------

To demonstrate the conditioning capability directly, we take exact
classical-shadow snapshots of the $4\times4$ transverse-field Ising torus,
hide the outcomes of half of the sites, and let Diffushadow complete them.
The completions are scored on cross-partition correlations
$\langle Z_i Z_j\rangle$ between an observed site $i$ and a completed site
$j$, for which the exact values are known. For both a contiguous mask (which
a fixed-order autoregressive model could also condition on) and an
interleaved mask (which it cannot), the inpainted cross-partition
correlations reproduce the exact values across all phases --- e.g.\
$1.003$ vs $1.009$ at $g=0.2$ and $0.289$ vs $0.296$ near criticality for
the interleaved mask --- while completions drawn without conditioning
collapse to near zero ($|\cdot|\le 0.12$), confirming that the model
performs genuine conditional inference on partially observed measurement
records. This operation --- completing an arbitrarily masked measurement
snapshot --- has no analogue in Hamiltonian-based numerical methods and is
unavailable to fixed-order autoregressive factorizations.

% Extended Data table (fill from inpaint_demo.npz):
%   g     mask         exact    inpainted   unconditional
%   0.20  prefix       +1.006   +0.997      +0.122
%   0.20  interleaved  +1.009   +1.003      +0.017
%   0.50  prefix       +0.935   +0.942      +0.012
%   0.50  interleaved  +0.901   +0.915      +0.052
%   0.75  prefix       +0.262   +0.240      +0.035
%   0.75  interleaved  +0.296   +0.289      -0.003
%   0.90  prefix       +0.075   +0.069      +0.024
%   0.90  interleaved  +0.041   +0.015      -0.020

% ---------------------------------------------------------------
% (4) Note for the Discussion limitations paragraph (one sentence):
% ---------------------------------------------------------------

On the benchmarks studied here, an equally sized autoregressive model
attains comparable cross-scale accuracy on fixed site-averaged
observables; the practical case for masked diffusion rests on its
conditioning flexibility and adaptive decoding rather than on accuracy
alone, and identifying observable classes or Hamiltonian families where
the factorizations genuinely separate remains an open question.
